% Source pp. 144-160 (Lectures13-18.pdf).
\section[Perturbative Renormalization Group]{Perturbative Renormalization Group\protect\footnote{PS Sect's 12.2, 12.3.}}

\subsection*{Finite renormalizations}

Renormalized perturbation theory allows one to obtain the correlation functions
of $\varphi^4$ theory as power series in the renormalized coupling parameter
$\lambda$ with well defined finite coefficients. This concerns the correlation
functions of the field $\varphi$ as well as the ones involving appropriately
renormalized composite fields. This is achieved by adjusting the counterterm
coefficients $\delta m^2$, $\delta\lambda$, $\delta Z$ in the action (and,
additionally, the renormalization factors in the composite fields) as the
functions of regularization parameter ($\epsilon=4-d$ in dimensional
regularization, or $\Lambda$ if some cutoff regularization is assumed).
Renormalized correlation functions appearing in this way depend on two
parameters, $\lambda$ and $m^2$, named the ``renormalized coupling'' and the
``renormalized mass parameter''. Precise relation of these parameters to
physical quantities, as well as the field normalizations, are specified
(explicitly or implicitly) by the \emph{renormalization scheme}.

By doing calculations in two different renormalization schemes we are not
solving two different field theories; rather, we are describing the same theory
in terms of different parameters. One can think of two-dimensional (because two
parameters, $m^2$ and $\lambda$, are involved) ``manifold of $\varphi^4$
theories'', with the parameters $m^2$ and $\lambda$ playing the role of (local)
coordinates on this manifold (the latter is often called the ``theory space'').

Each point of this manifold represents a field theory, i.e.\ a collection of all
correlation functions, defined up to finite field renormalizations
\[
\bigl\langle\varphi(x_1)\cdots\varphi(x_n)\bigr\rangle\ \sim\ Z_{\text{finite}}^{n/2}\,\bigl\langle\varphi(x_1)\cdots\varphi(x_n)\bigr\rangle\,.
\]
To simplify things one can consider the ratios like
\[
R(x_1,x_2;x_3,\dots,x_n)=\bigl\langle\varphi(x_1)\varphi(x_2)\cdots\varphi(x_{2n})\bigr\rangle\big/\bigl\langle\varphi(x_1)\varphi(x_2)\bigr\rangle^n
\]
which do not depend on the normalizations of the fields. Different
renormalization schemes correspond to different choices of the coordinate system
on the above manifold. Once the coordinate system (i.e.\ renormalization scheme)
is fixed, that is the points are somehow ``labeled'' by two parameters $m^2$ and
$\lambda$, the above ratios become functions of these parameters,
\[
R(x_i|m^2,\lambda)\,.
\]
Relation between the parameters $(m^2,\lambda)$ and $(\tilde m^2,\tilde\lambda)$,
associated with two different renormalization schemes, is just a coordinate
transformation,
\[
m^2\to\tilde m^2(m^2,\lambda)\,,\qquad\lambda\to\tilde\lambda(m^2,\lambda)\,.
\]
At $d=4$ this relation has the form (by dimensional analysis)
\begin{equation}\label{eq:schemechange}
\tilde m^2=m^2\,f(\lambda)\,,\qquad\tilde\lambda=g(\lambda)\,,
\end{equation}
where the functions $f,g$ can in principle be computed order by order in
$\lambda$,
\begin{align*}
f(\lambda)&=f_0+\lambda\,f_1+\lambda^2\,f_2+\cdots\,,\\
g(\lambda)&=\lambda+\lambda^2\,g_2+\lambda^3\,g_3+\cdots\,,
\end{align*}
the coefficients $f_k$, $g_k$ being finite numbers. The correlation functions
(or rather the above ratios $R$) computed in the two schemes are related as
\[
\tilde R\bigl(x_i|m^2f(\lambda),g(\lambda)\bigr)=R(x_i|m^2,\lambda)\,,
\]
so that these are the same correlation functions expressed in terms of different
parameters, that is, coordinates in the space of $\varphi^4$ theories.

\subsection*{Short distance problem}

Renormalized perturbation theory, no matter what scheme is assumed, yields the
correlation functions as formal power series in the coupling parameter
$\lambda$. Although in each given order these series give finite and well
behaved correlation functions, the problem of overall \emph{non-perturbative}
consistency of the theory is a separate question. There is an issue of
convergence of the perturbative expansion. It is possible to show that the
series in $\varphi^4$ have in fact zero radius of convergence, and at the best
they can be understood as asymptotic series. If so, there is an important
question if these series can be understood as asymptotic expansions of some
non-perturbative correlation functions which are compatible with requirements of
\emph{locality} and with \emph{quantum mechanical interpretation} (the property
usually called \emph{positivity}). Although one can show that these properties
(locality and positivity) are preserved at the perturbative level, i.e.\ in each
given order in $\lambda$, it is conceivable that they can break down after
summing the series up.

We have seen that in Euclidean version of quantum field theory locality implies
in particular that the correlation function
\[
\bigl\langle\varphi(x_1)\varphi(x_2)\cdots\varphi(x_{2n})\bigr\rangle\,,
\]
viewed as the function of, say, $x_1$, has only point-like singularities at
$x_1=x_2,x_3,\dots,x_{2n}$. (We have observed this in the free Klein-Gordon
theory (see Section~4) for the two-point function, but the arguments we have
applied in that analysis depend only on general properties of the
energy-momentum spectrum, and are generalized straightforwardly to general
interacting theories; I hope to come back to this point later.) This is because
upon continuation to complex coordinates and ultimately to the Minkowski
space-time these singularities translate into the singularities at the light
cones associated with the points $x_2,\dots,x_{2n}$; we have seen that these
singularities are directly related to commutators of the field operators (see
again the analysis in Section~4). Locality states that the commutators of
Heisenberg field operators must vanish outside the light cone,
\[
[\hat\varphi(\xx,t),\hat\varphi(\mathbf0,0)]=0\qquad\text{for}\qquad t^2<\xx^2\,.
\]
This property of locality is preserved in the renormalized perturbation theory at
each order in $\lambda$ but there is no guarantee it will not break down in full
non-perturbative theory. It is very well conceivable that upon summing the
perturbative series up singularities at finite space-like separations may be
generated, leading to violations of the local commutativity.

The following example demonstrates this possibility. Consider the power series
\[
\sum_{n=0}^\infty(-\lambda)^n\,\Bigl(\log m|x_1-x_2|\Bigr)^n\,,
\]
which models a perturbative series for some correlation function. At each given
order the above expression has point-like singularity at $x_1=x_2$. However its
sum is
\[
\frac1{1+\lambda\log m|x_1-x_2|}\,.
\]
This function is singular at
\[
|x_1-x_2|=x_0\equiv m^{-1}\,e^{-1/\lambda}\,.
\]
As the result, if one continues to the Minkowski space-time the light cone
``fattens'' and causality is broken at short scales $\simeq x_0$.\footnote{One
can take different attitudes toward such possibility. One pragmatic attitude is
that we do not know what describes short-distance physics anyway. After all, at
short distances quantum gravity must play important role, perhaps along with
extra dimensions, superstrings, D-branes, and whatnot. And therefore we do not
care about possible inconsistencies at these short distances. With this attitude
quantum field theory is viewed as a way of producing (renormalized) perturbative
expansions in the coupling parameter. Such expansions of course are useful only
if the coupling parameter is small. In quantum electrodynamics (QED) where the
analog of $\lambda$ is the fine structure constant which is small indeed, this is
valid point of view. And if the interaction is not small in any sense, according
to this philosophy one would conclude that quantum field theory is not useful at
all.\par
In the past, at certain stage of development, such was prevailing attitude.
Discovery of \emph{Asymptotic Freedom} in nonabelian gauge theories, along with
better understanding of \emph{Renormalization Group} has changed that. On the
face of it, the discovery of the asymptotic freedom was pure perturbative result,
obtained by computing few diagrams, but through understanding the Renormalization
Group one can see that it actually guaranteed full nonperturbative consistency of
quantum nonabelian gauge theories. More generally, renormalization allows one to
understand quantum field theory in terms of criticality, or vice versa, to the
extent that from theoretical point of view quantum field theory and theory of
critical phenomena are much one and the same thing. Moreover, main idea of the
Wilson's renormalization group has proved very useful in theoretical physics well
beyond quantum field theory. This is why I hope to devote significant part of this
course to it.}

The problem of the short-distance consistency of a quantum field theory requires
analysis of its behavior under changes of the length scale. The latter is
generally understood in terms of \emph{Renormalization Group}. Here I will
briefly discuss the basics of this approach, in its perturbative version, and
applied to $\varphi^4$ theory. The analysis is substantially simpler in the case
of \emph{massless theory}.

\subsection*{Massless \texorpdfstring{$\varphi^4$}{phi**4}}

As a preparation, let us discuss what happens with our renormalized perturbation
if one sets the renormalized mass parameter $m^2$ to zero. Although precise
meaning of $m^2$ depends on the renormalization scheme, the condition
\[
m^2=0
\]
is \emph{scheme independent}, i.e.\ once it is satisfied in one scheme, it is
satisfied in any other scheme, as the Eq.~\eqref{eq:schemechange} suggests.
There are at least two reasons why this case is interesting.
\begin{itemize}
\item Potential problem with locality can occur at short scales, much shorter
than $m^{-1}$. The mass term is not important at such scales, and by analyzing
this case one can understand the general situation.
\item We will see that in applications to critical phenomena the case $m^2=0$
describes physics at exactly critical point, and therefore it is of central
interest.
\end{itemize}

Taking a look at our previous perturbative results we discover that it is not
possible to just set $m^2=0$ there, since nearly all quantities we have computed
are singular at this point. Consider for instance the 4-point vertex,
\[
\Gamma^{(4)}(p_1,\dots,p_4)=\lambda-\frac{\lambda^2}2\Bigl(I_r(p_{12}^2)+I_r(p_{13}^2)+I_r(p_{14}^2)\Bigr)+O(\lambda^3)\,,
\]
where again
\[
I_r(p^2)=-\frac1{(4\pi)^2}\int_0^1\dd u\,\log\left(1+u(1-u)\,\frac{p^2}{m^2}\right)\,.
\]
Obviously, there is a singularity at $m^2=0$. This can signal one of two things.
Either the theory itself becomes sick at this point, or our coordinate system
does not extend smoothly to this case. I want to argue that $m^2=0$
singularities of the perturbative correlation functions originate from a
coordinate singularity.\footnote{Think of some a function $f(x,y)$ written in new
coordinates $(x',y')$, $x=x'$, $y=y'+y'^2\log x'^2$.}

Indeed, a particular coordinate system is determined by the normalization
conditions. For example, in our scheme SII we have assumed
\begin{equation}\label{eq:SIIrecall}
\Gamma^{(2)}(p^2)\Big|_{p^2=0}=m^2\,,\qquad\text{and}\qquad\frac{\dd\Gamma^{(2)}(p^2)}{\dd p^2}\Big|_{p^2=0}=1\,.
\end{equation}
At $m^2=0$ the first of these conditions remains useful, just stating that
$\Gamma^{(2)}(p^2)$ turns to zero at $p^2=0$, which can be considered as a
scheme-independent definition of the ``massless theory'', but at $m^2=0$ this is
no longer a simple zero, being replaced by more complicated singular point. We
will see in a moment that at $m^2=0$ $\Gamma^{(2)}$ contains terms like
$p^2\log p^2$ and therefore the second of the conditions \eqref{eq:SIIrecall}
cannot be imposed. The coordinate system associated with our scheme is singular
at $m^2=0$.

It is useful to have a scheme which admits smooth limit $m^2\to0$. Such scheme
must involve dimensional parameter different from $m^2$. Let us denote it
$\mu^2$, $[\mu]=1$. Since $\mu$ can be chosen arbitrarily, we will have families
of schemes, parameterized by auxiliary scale $\mu$. Let us take one such
$\mu$-dependent scheme, defined by the normalization conditions
\begin{align}
\frac{\dd\Gamma^{(2)}(p^2)}{\dd p^2}\Big|_{p^2=\mu^2}&=1\,,\label{eq:mu1}\\
\Gamma^{(4)}(p_1,\dots,p_4)\Big|_{p_i=\mu e_i}&=\lambda_{(\mu)}\,,\label{eq:mu2}
\end{align}
where $e_i$ is some set of four fixed 4-vectors, which satisfy
$e_1+e_2+e_3+e_4=0$, but have no other linear dependencies. To be specific, we
can choose symmetric arrangement,
\[
e_i^2=3\qquad\text{and}\qquad e_i\cdot e_j=-1\quad\text{for}\quad i\neq j\,.
\]
I write $\lambda_{(\mu)}$ for the renormalized coupling constant as the reminder
about the scale $\mu$ associated with this scheme. It is conventional to refer to
this auxiliary scale as the ``normalization scale''. It is perfectly legitimate
to use such scheme in general case, but we will apply it to massless theory,
which is singled out by the condition
\begin{equation}\label{eq:mu3}
\Gamma^{(2)}(0)=0\,.
\end{equation}

We can now develop renormalized perturbation theory for massless $\varphi^4$
theory taking exact same steps we did before. We start with the action in its
bare form,
\[
\mathcal A=\int \dd^dx\,\left[\frac12\,(\partial\varphi_0)^2+\frac{m_0^2}2\,\varphi_0^2+\frac{\lambda_0}{4!}\,\varphi_0^4\right]\,,
\]
with some bare mass parameter (yet to be adjusted), and write it as
\[
\mathcal A=\int \dd^dx\,\left[\frac12\,(\partial\varphi)^2+\frac{\lambda_{(\mu)}}{4!}\,\varphi^4+\text{counterterms}\right]\,,
\]
where
\[
\varphi=Z_{(\mu)}^{-1/2}\,\varphi_0\,,
\]
and the counterterms are
\[
\text{counterterms}=\frac{\delta Z_{(\mu)}}2\,(\partial\varphi)^2+\frac{\delta m^2_{(\mu)}}2\,\varphi^2+\frac{\delta\lambda_{(\mu)}}{4!}\,\varphi^4\,,
\]
with
\[
\delta Z_{(\mu)}=Z_{(\mu)}-1\,,\qquad
\delta m^2_{(\mu)}=m_0^2\,Z_{(\mu)}\,,\qquad
\delta\lambda_{(\mu)}=\lambda_0\,Z_{(\mu)}^2-\lambda_{(\mu)}\,.
\]
Again, we assume that all the counterterm coefficients are power series in the
renormalized coupling $\lambda_{(\mu)}$,
\begin{align*}
\delta Z_{(\mu)}&=z^{(2)}_{(\mu)}\,\lambda_{(\mu)}^2+\cdots\,,\\
\delta\lambda_{(\mu)}&=a^{(2)}_{(\mu)}\,\lambda_{(\mu)}^2+\cdots\,,\\
\delta m^2_{(\mu)}&=b^{(1)}_{(\mu)}\,\lambda_{(\mu)}+\cdots\,,
\end{align*}
with the coefficients determined from the normalization conditions
\eqref{eq:mu1}, \eqref{eq:mu2}, \eqref{eq:mu3}.

A remark is in order. We have put two parameters, $m_0^2$ and $\lambda_0$, in the
bare action. Renormalized massless theory depends on a single parameter, the
renormalized coupling $\lambda_{(\mu)}$, only. The normalization scale $\mu$ is
not a parameter of the theory, in that it is an attribute of the renormalization
scheme, rather than the theory itself. We have imposed three normalization
conditions Eqs.~\eqref{eq:mu1}--\eqref{eq:mu3}. While Eq.~\eqref{eq:mu2} defines
the renormalized coupling $\lambda_{(\mu)}$ and Eq.~\eqref{eq:mu1} fixes
normalization of $\varphi$, the Eq.~\eqref{eq:mu3} states that the theory is
massless, and thus it imposes restriction on the theory itself. In principle,
this last condition can be translated into certain relation between the bare
parameters $m_0^2$ and $\lambda_0$. The bare theory is implicitly equipped with
some regularization, and precise form of such relation depends on the
regularization method. For instance, if some cutoff regularization is implied,
with the cutoff momentum $\Lambda$, at $d=4$ the relation has general form
\begin{equation}\label{eq:critbare}
m_0^2=\Lambda^2\,C(\lambda_0)\,.
\end{equation}
Precise shape of $C(\lambda_0)$ heavily depends on specific implementation of the
cutoff. Subset of $\varphi^4$ theories selected by this condition corresponds to
critical domain in the theory of second order phase transitions. In dimensional
regularization there is no cutoff parameter $\Lambda$, and simple dimensional
analysis shows that this relation must take simple form
\[
m_0^2=0\qquad\text{(dimensional regularization)}.
\]
This means in particular that $\delta m^2=0$, and we do not have to bother with
the mass renormalization at all. This is one of the technical advantages of the
dimensional regularization technique.

Again, thinking in terms of the ``two-dimensional manifold of $\varphi^4$
theories'', the massless $\varphi^4$ theories correspond to certain
one-dimensional ``critical'' submanifold (critical curve, in our case),
figure~\ref{fig:critcurve}.
\begin{figure}[ht]
\centering\figfile{critcurve}
\caption{The two-dimensional space of $\varphi^4$ theories, with coordinates
$(m^2,\lambda)$. The massless theories form the critical curve $m^2=0$.}
\label{fig:critcurve}
\end{figure}
In statistical mechanics, the ``bare action'' models microscopic interaction, and
the parameters $m_0^2$ and $\lambda_0$, as well as the cutoff energy $\Lambda$,
are related to the temperature and the microscopic physical characteristics of
the matter under consideration (e.g.\ molecular interactions in a gas/liquid).
When these parameters cross the ``critical submanifold'' the macroscopic system
undergoes a second order phase transition.

\subsection*{Renormalization as the scale transformation. Callan-Symanzik equation}

Although in general $\varphi^4$ theory the coordinates $(m^2,\lambda)$ depend on
the renormalization scheme, the ``critical submanifold'' $m^2=0$ is scheme
independent. It is because the condition \eqref{eq:mu3}, being a physical
statement, is scheme independent. But even within the critical curve there still
is a freedom in the choice of the coordinate $\lambda$ along the curve. Within
the above class of schemes \eqref{eq:mu1}--\eqref{eq:mu2} this choice is
controlled by the choice of the normalization scale $\mu$. If $\mu$ and
$\tilde\mu$ are two different values of the normalization scale, the associated
renormalized couplings constants
\[
\lambda=\lambda_{(\mu)}\qquad\text{and}\qquad\tilde\lambda=\lambda_{(\tilde\mu)}
\]
are related as
\[
\tilde\lambda=g(\lambda)=g(\lambda|\tilde\mu/\mu)\,,
\]
where at $d=4$, by dimensional counting, the function $g$ can depend on
$\lambda$ and the dimensionless ratio $\tilde\mu/\mu$. And if $\tilde\mu$ is
infinitesimally close to $\mu$, i.e.\ $\tilde\mu=\mu+\delta\mu$, then
\[
\tilde\lambda=\lambda+\frac{\delta\mu}\mu\,\beta(\lambda)\,,\qquad\text{where}\qquad
\beta(\lambda)=\frac{\dd g(\lambda|\xi)}{\dd\xi}\Big|_{\xi=1}
\]
is a function of $\lambda$ alone.

Of course the correlation functions are essentially independent of the
renormalization scheme, in the present context of the choice of the scale $\mu$.
More precisely, the ratios of the correlation functions, e.g.
\begin{equation}\label{eq:Rratio}
R(x_i|\lambda_{(\mu)},\mu)=\frac{\bigl\langle\varphi(x_1)\cdots\varphi(x_{2n})\bigr\rangle_{(\mu)}}{\langle\varphi\varphi\rangle^n_{(\mu)}}
\end{equation}
which are insensitive to the normalization of $\varphi$, do not change under
substitutions
\[
\mu\to\tilde\mu\,,\qquad\lambda_{(\mu)}\to\lambda_{(\tilde\mu)}\,.
\]
Here the index $\mu$ in $\langle\cdots\rangle_{(\mu)}$ signifies that the
correlation function is calculated in the scheme with the normalization scale
$\mu$. This independence can be expressed in the form of partial differential
equation
\[
\left(\mu\frac\partial{\partial\mu}+\beta(\lambda)\,\frac\partial{\partial\lambda}\right)R(x_i|\lambda,\mu)=0\,,
\]
where now $\lambda$ and $\mu$ are viewed as two independent variables.

These equations can be extended to the correlation functions themselves if one
recalls that a change in renormalization scheme in general induces a finite
change in normalization of $\varphi$,
\[
\varphi_{(\tilde\mu)}=Z^{-1/2}_{\text{finite}}(\lambda|\tilde\mu/\mu)\,\varphi_{(\mu)}\,,
\]
where again at $d=4$ the constant $Z_{\text{finite}}$ can depend on $\lambda$ and
dimensionless ratio $\tilde\mu/\mu$. For an infinitesimal change
$\tilde\mu=\mu+\delta\mu$
\[
\varphi_{(\mu+\delta\mu)}=\left(1-\frac{\delta\mu}\mu\,\gamma(\lambda)\right)\,\varphi_{(\mu)}\,,
\qquad\text{with}\qquad
\gamma(\lambda)=\frac12\,\frac{\dd Z_{\text{finite}}(\lambda,\xi)}{\dd\xi}\Big|_{\xi=1}\,.
\]
Taking into account this transformation of $\varphi$ we obtain the so-called
\emph{Callan-Symanzik equation}:
\begin{equation}\label{eq:CS}
\left(\mu\frac\partial{\partial\mu}+\beta(\lambda)\,\frac\partial{\partial\lambda}+n\,\gamma(\lambda)\right)\bigl\langle\varphi(x_1)\cdots\varphi(x_n)\bigr\rangle_{(\mu)}=0
\end{equation}
where the correlation function is viewed as depending on two parameters,
$\lambda$ and $\mu$.

At the first glance the Callan-Symanzik equation contains strictly zero physical
information. Indeed, in order to make a meaningful renormalized perturbation
theory we have introduced an auxiliary scale $\mu$ having in mind that nothing
physical can depend on this arbitrary scale, and then the equation \eqref{eq:CS}
just states that yes indeed nothing physical depends on $\mu$. This is
absolutely true if no information about the functions $\beta(\lambda)$ and
$\gamma(\lambda)$ is added. However, if something is known about those functions
the Callan-Symanzik equation becomes useful.

To see this let us first notice that by dimensional counting any correlation
function of the $d=4$ massless theory can be written as
\[
\bigl\langle\varphi(x_1)\cdots\varphi(x_n)\bigr\rangle_{(\mu)}=\mu^n\,C^{(n)}(\mu x_1,\dots,\mu x_n|\lambda)\,,
\]
where $C^{(n)}$ is dimensionless function of dimensionless variables $\mu x_i$
and $\lambda$. Therefore the Callan-Symanzik equation \eqref{eq:CS} can be written
as
\begin{equation}\label{eq:CSdil}
\left(\sum_{i=1}^nx_i\frac\partial{\partial x_i}+n\bigl(1+\gamma(\lambda)\bigr)+\beta(\lambda)\,\frac\partial{\partial\lambda}\right)\bigl\langle\varphi(x_1)\cdots\varphi(x_n)\bigr\rangle=0\,,
\end{equation}
where the coordinate derivative operators have the following exact meaning
\[
x_i\frac\partial{\partial x_i}=\sum_{\mu=1}^4x_i^\mu\,\frac\partial{\partial x_i^\mu}\,.
\]
Note that the differential operator $\sum_ix_i\partial/\partial x_i$ describes
infinitesimal dilation of the coordinates $x_i$ and therefore the equation
\eqref{eq:CSdil} is the statement that such dilation can be essentially
compensated (up to the field renormalization) by appropriate change of the
coupling parameter $\lambda$. More precisely, an infinitesimal dilation of the
coordinates is equivalent to certain infinitesimal shift of $\lambda$ plus an
infinitesimal renormalization of $\varphi$,
\[
x^\mu\to(1+\delta L)\,x^\mu\quad\Longleftrightarrow\quad
\lambda\to\lambda-\beta(\lambda)\,\delta L\,,\quad\varphi\to\varphi-\bigl(1+\gamma(\lambda)\bigr)\,\varphi\,\delta L\,.
\]
Let $\lambda(L)$ be a solution of the following differential equation and initial
condition
\begin{equation}\label{eq:flow1}
L\frac{\dd}{\dd L}\,\lambda(L)=-\beta\bigl(\lambda(L)\bigr)\,,\qquad\lambda(1)=\lambda\,.
\end{equation}
Then the Callan-Symanzik equation states that
\begin{equation}\label{eq:CSsol}
Z^{n/2}(L)\,\bigl\langle\varphi(x_1/L)\cdots\varphi(x_n/L)\bigr\rangle_{\lambda(L)}=\bigl\langle\varphi(x_1)\cdots\varphi(x_n)\bigr\rangle_\lambda\,,
\end{equation}
where the index $\lambda$ now shows the value of renormalized coupling constant,
and $Z(L)$ solves the linear differential equation
\begin{equation}\label{eq:ZL}
L\frac{\dd}{\dd L}\,Z^{1/2}(L)=-\Bigl(1+\gamma\bigl(\lambda(L)\bigr)\Bigr)\,Z^{1/2}(L)\,,\qquad Z^{1/2}(1)=1\,.
\end{equation}
It is assumed that the correlation functions in \eqref{eq:CSsol} are computed
within the same scheme, with the same normalization scale $\mu$, but with
different values, $\lambda(L)$ and $\lambda$, of the renormalized coupling
parameter. Note that $\mu$ does not appear in \eqref{eq:CSsol}, i.e.\ in this form
the Callan-Symanzik equation is not a statement about scheme dependence but
rather determines how the theory behaves when one goes from one scale to
another.

\subsection*{Renormalization group flow}

We can give this equation the following geometric interpretation. The coupling
constants $\lambda$ and $\lambda(L)$ are coordinates representing two points, $P$
and $P(L)$, of the critical submanifold (figure~\ref{fig:flowpoints}).
\begin{figure}[ht]
\centering\figfile{flowpoints}
\caption{Two points $P$ and $P(L)$ on the critical curve, with coordinates
$\lambda$ and $\lambda(L)$.}
\label{fig:flowpoints}
\end{figure}
The equation \eqref{eq:CSsol} then states that the two field theories $P$ and
$P(L)$ are related by the scale transformation
\begin{equation}\label{eq:scaletr}
x\to L\,x
\end{equation}
(the field renormalizations obviously do not affect the physical content of the
theory). The equation
\begin{equation}\label{eq:RGflow}
L\frac{\dd}{\dd L}\,\lambda(L)=-\beta\bigl(\lambda(L)\bigr)
\end{equation}
describes, in given coordinate system, the ``flow'' in the space of field
theories generated by the scale transformations. It is called the
\emph{renormalization group flow}, and the equation \eqref{eq:RGflow} is usually
referred to as the \emph{renormalization group flow equation}. The quantity
$\lambda(L)$ is often referred to as ``running'', or ``scale-dependent'' coupling
constant. The fact that the scale transformations \eqref{eq:scaletr} generate a
``flow'' is not limited to the massless theory; we will see shortly that the
Callan-Symanzik equation generalizes to the case $m^2\neq0$.

We have seen that the functions $\beta(\lambda)$ and $\gamma(\lambda)$ contain
most important information about the scale dependence of the massless theory.
Unfortunately, these functions are not known exactly. However we can compute them
perturbatively, order by order in the coupling parameter $\lambda$.

There are many ways to do this. The simplest is just to compute the correlation
functions perturbatively and then determine the functions $\beta$ and $\gamma$
from the Callan-Symanzik equation \eqref{eq:CS} itself. Let us first note that the
Callan-Symanzik equation in the form analogous to \eqref{eq:CS} applies to the
momentum-space correlation functions as well as to coordinate-space ones (the
Fourier transform relating the two depends neither on $\lambda$ nor on $\mu$).
Moreover, simple calculation shows that this equation can be written as the
equation
\begin{equation}\label{eq:CSGamma}
\left(\mu\frac\partial{\partial\mu}+\beta(\lambda)\,\frac\partial{\partial\lambda}-n\,\gamma(\lambda)\right)\Gamma^{(n)}(p_1,\dots,p_n)=0
\end{equation}
involving the proper vertices.

To the order $\lambda^2$ the proper vertices $\Gamma^{(2)}$ and $\Gamma^{(4)}$,
when computed in the scheme \eqref{eq:mu1}--\eqref{eq:mu2}, are (I ask you to
carry out the calculations as the homework Problem~\ref{hw:5})
\begin{align}
\Gamma^{(2)}(p^2)&=p^2+\frac1{12}\,\frac{\lambda^2}{(4\pi)^4}\,p^2\,\Bigl(\log(\mu^2/p^2)+1\Bigr)+O(\lambda^3)\,,\label{eq:G2mu}\\
\Gamma^{(4)}(p_1,\dots,p_4)&=\lambda+\frac{\lambda^2}{2\,(4\pi)^2}\left[\log\left(\frac{p_{12}^2}{4\mu^2}\right)+\log\left(\frac{p_{13}^2}{4\mu^2}\right)+\log\left(\frac{p_{14}^2}{4\mu^2}\right)\right]+O(\lambda^3)\,,\label{eq:G4mu}
\end{align}
where $p_{12}^2=(p_1+p_2)^2$, etc. Let us first plug \eqref{eq:G2mu} into
\eqref{eq:CSGamma}. The result is
\[
\frac2{12}\,\frac{\lambda^2}{(4\pi)^4}\,p^2+\beta(\lambda)\,\frac2{12}\,\frac\lambda{(4\pi)^4}\,p^2\,\bigl[\log(\mu^2/p^2)+1\bigr]-2\gamma(\lambda)\,p^2+O(\lambda^3)=0\,,
\]
where the first term came from the $\mu\partial/\partial\mu$ term in
\eqref{eq:CSGamma}. The second term contains $p^2\log\mu^2/p^2$; this dependence
cannot be canceled by two other terms exposed in the above equation. The only way
to satisfy this equation in the order $\lambda^2$ is to assume that
$\beta(\lambda)\sim\lambda^2$ so that the second term has to be attributed to the
$O(\lambda^3)$ corrections. We then find
\begin{equation}\label{eq:gamma1}
\gamma(\lambda)=\frac1{12}\,\frac{\lambda^2}{(4\pi)^4}+O(\lambda^3)\,.
\end{equation}
Similarly, substitute the above perturbative result for $\Gamma^{(4)}$ into
\eqref{eq:CSGamma}. We have
\[
-3\,\frac{\lambda^2}{(4\pi)^2}+\beta(\lambda)-4\lambda\,\gamma(\lambda)+O(\lambda^3)=0\,.
\]
We already know that $\gamma(\lambda)\sim\lambda^2$ and so the term involving this
function is $O(\lambda^3)$. Therefore
\begin{equation}\label{eq:beta1}
\beta(\lambda)=3\,\frac{\lambda^2}{(4\pi)^2}+O(\lambda^3)\,.
\end{equation}

Note that the leading term \eqref{eq:beta1} of the $\beta$-function is
\emph{positive}. This is the feature of $\varphi^4$ theory. The flow equation
\eqref{eq:flow1} then shows that for small $\lambda$ the scale-dependent coupling
constant $\lambda(L)$ \emph{decreases} with $L$. The equation \eqref{eq:CSsol}
shows that larger $L$ correspond to \emph{longer} length scales (this is clear if
one rewrites \eqref{eq:CSsol} in equivalent form
\[
\bigl\langle\varphi(x_1)\cdots\varphi(x_n)\bigr\rangle_{\lambda(L)}\sim\bigl\langle\varphi(Lx_1)\cdots\varphi(Lx_n)\bigr\rangle_\lambda
\]
showing that the correlation functions computed with $\lambda(L)$ are equivalent
to the correlation functions computed with $\lambda$, with the separations
``stretched'' by the scale factor $L$).

If one ignores the higher order corrections in \eqref{eq:beta1} the flow equation
\eqref{eq:flow1} is easily integrated,
\begin{equation}\label{eq:running}
\lambda(L)=\frac\lambda{1+3\frac\lambda{(4\pi)^2}\log L}\,.
\end{equation}
At sufficiently small $L$ (i.e.\ sufficiently short distances) the function
$\lambda(L)$ diverges. This signals potential trouble with locality which we
discussed earlier. Of course when $\lambda(L)$ becomes large one cannot ignore the
higher-order corrections in \eqref{eq:beta1}; in fact in order to understand
short-distance behavior of $\varphi^4$ theory we need to know the shape of the
$\beta$-function at large $\lambda$. We see that the question of consistency of
$\varphi^4$ theory cannot be settled on the perturbative level.

In nonabelian gauge theories the leading term of the $\beta$-function turns out
to be \emph{negative}, meaning that the running coupling constant
\emph{decreases} at short distances (instead it increases at large distances); as
a result, although nonabelian gauge theory usually cannot be solved by means of
perturbation theory, the question of its \emph{consistency} as a local field
theory can be positively answered on the basis of perturbative calculations. We
will discuss this phenomenon known as \emph{asymptotic freedom} in greater details
later.

\subsection*{Mass perturbation}

Let us refine our techniques further. So far we have considered the
renormalization group flow in massless $\varphi^4$ theory. We would like to
generalize this picture to the massive case. The most natural way to do this is
to consider the massive $\varphi^4$ theory as the massless theory perturbed by
adding the term $\frac{m^2}2\varphi^2$ to the action density. We already know that
the field $\varphi^2$ requires renormalization. Let us discuss how this field
renormalizes in our $\mu$-dependent scheme.

In the massless theory the normalization condition
$\Gamma^{(\varphi^2_R,2)}(0,0)=1$ we have used before has to be replaced by some
nonsingular $\mu$ dependent condition. One can take for example
\begin{equation}\label{eq:phi2mu}
\Gamma^{(\varphi^2_R,2)}(p_1,p_2)\Big|_{\substack{p_i^2=2\mu^2\\p_1\cdot p_2=-\mu^2}}=1\,.
\end{equation}
The renormalized field $[\varphi^2]_{(\mu)}$ defined by this condition differs
from the bare field $\varphi_0^2$ by the renormalization
constant\footnote{Note that this definition of the $\varphi^2$ renormalization
constant differs from that introduced in Section~14, Eq.~\eqref{eq:phi2Rinv}.
There, we have defined the constant $Z_{\varphi^2}$ as the coefficient between
$[\varphi^2]_R$ and $\varphi^2$. The new $Z^{\varphi^2}$ differs from
$Z_{\varphi^2}$ in Section~14 by the factor of $Z$.},
\[
[\varphi^2]_{(\mu)}=\bigl(Z^{\varphi^2}_{(\mu)}\bigr)^{-1}\,\varphi_0^2\,.
\]
Simple calculation similar to that discussed in the previous Section shows that
in the leading order
\begin{equation}\label{eq:Zphi2mu}
Z^{\varphi^2}_{(\mu)}=1-\frac12\,\frac\lambda{(4\pi)^{d/2}}\,(2\mu^2)^{d/2-2}\,\frac{\Gamma(2-d/2)\,\Gamma^2(d/2-1)}{\Gamma(d-2)}+O(\lambda^2)
\end{equation}
Although $Z^{\varphi^2}_{(\mu)}$ is divergent at $d=4$, the ratio
$Z^{\varphi^2}_{(\tilde\mu)}/Z^{\varphi^2}_{(\mu)}$ is expected to be finite. An
infinitesimal change of the renormalization scale $\mu\to\mu+\delta\mu$ then leads
to infinitesimal renormalization
\[
[\varphi^2]_{(\mu+\delta\mu)}=\left(1-\frac{\delta\mu}\mu\,\gamma_{\varphi^2}(\lambda)\right)\,[\varphi^2]_{(\mu)}\,,
\]
where
\begin{equation}\label{eq:gammaphi2}
\gamma_{\varphi^2}(\lambda)=\bigl(Z^{\varphi^2}_{(\mu)}\bigr)^{-1}\,\mu\frac{\dd}{\dd\mu}\,Z^{\varphi^2}_{(\mu)}
\end{equation}
is finite function of $\lambda$. From \eqref{eq:Zphi2mu} we obtain
\begin{equation}\label{eq:gammaphi2val}
\gamma_{\varphi^2}(\lambda)=\frac\lambda{(4\pi)^2}+O(\lambda^2)\,.
\end{equation}
Repeating our previous calculations we find that a correlation function with $N$
insertions of $[\varphi^2]_{(\mu)}$ and $n$ insertions of $\varphi$, satisfies the
modified Callan-Symanzik equation
\[
\left(\mu\frac\partial{\partial\mu}+\beta(\lambda)\,\frac\partial{\partial\lambda}+n\,\gamma(\lambda)+N\,\gamma_{\varphi^2}(\lambda)\right)
\bigl\langle\varphi(x_1)\cdots\varphi(x_n)\,[\varphi^2]_{(\mu)}(y_1)\cdots[\varphi^2]_{(\mu)}(y_N)\bigr\rangle_{(\mu)}=0\,.
\]

Let us now introduce the mass term as
\[
\mathcal A_{M^2}=\mathcal A_0+\int \dd^dx\,\frac{M^2}2\,[\varphi^2]_{(\mu)}\,,
\]
where $\mathcal A_0$ is the massless action
\[
\mathcal A_0=\int \dd^dx\,\left[\frac12\,(\partial\varphi)^2+\frac\lambda{4!}\,\varphi^4+\frac{\delta Z_{(\mu)}}2\,(\partial\varphi)^2+\frac{\delta\lambda_{(\mu)}}{4!}\,\varphi^4\right]\,.
\]
Here $\delta Z_{(\mu)}$ and $\delta\lambda_{(\mu)}$ are counterterms computed in the
massless theory, within the above $\mu$-dependent renormalization scheme (and I
assumed dimensional regularization and therefore the $\delta m^2$ counterterm is
not included). Recalling that
\[
[\varphi^2]_{(\mu)}=\frac{Z_{(\mu)}}{Z^{\varphi^2}_{(\mu)}}\,\varphi^2\,,\qquad Z_{(\mu)}=1+\delta Z_{(\mu)}\,,
\]
the massive action can be written as
\[
\mathcal A_{M^2}=\int \dd^dx\,\left[\frac12\,(\partial\varphi)^2+\frac{M^2}2\,\varphi^2+\frac\lambda{4!}\,\varphi^4+\frac{\delta Z_{(\mu)}}2\,(\partial\varphi)^2+\frac{\delta M^2_{(\mu)}}2\,\varphi^2+\frac{\delta\lambda_{(\mu)}}{4!}\,\varphi^4\right]\,,
\]
where
\begin{equation}\label{eq:deltaM2}
\delta M^2_{(\mu)}=M^2\left(\frac{Z_{(\mu)}}{Z^{\varphi^2}_{(\mu)}}-1\right)\,.
\end{equation}
This action can be used now to develop a renormalized perturbation theory in the
massive case. Note that the above definition of the mass parameter corresponds to
specific renormalization scheme in which
\begin{itemize}
\item The counterterms $\delta Z_{(\mu)}$ and $\delta\lambda_{(\mu)}$ are borrowed
from the massless theory; in particular these coefficients do not depend on $M^2$.
\item The counterterm $\delta M^2_{(\mu)}$ is expressed by \eqref{eq:deltaM2}
through the $\varphi^2$-renormalization constant $Z^{\varphi^2}_{(\mu)}$; again,
this constant has to be computed upfront in the massless theory, and
$\delta M^2_{(\mu)}$ depends on $M^2$ linearly.
\end{itemize}
Although it is possible to show that the counterterms thus defined indeed
compensate for the divergences in the renormalized perturbation theory, no
explicit normalization conditions are imposed in the massive theory. Relation
between $M^2$, $\lambda$ and physical quantities in this scheme is not specified
in advance but must be computed order by order in the perturbation theory. In
this sense this renormalization scheme is similar to the ``minimal subtraction''
scheme.

Note that in this scheme the correlation functions depend on two mass
parameters, $M$ and our auxiliary scale $\mu$, although we know only one is
needed. This is the price we pay for defining specific coordinate system
$(M^2,\lambda)$ which is nonsingular on the critical submanifold $M^2=0$. It is
not difficult to check that correlation functions computed within this scheme
admit smooth limit $M^2=0$; in fact it is clear in advance: at $M^2=0$ this
perturbation theory simply reduces to the massless renormalized perturbation
theory.

Like in the massless case, nothing except for normalizations depends on the
auxiliary scale $\mu$. Under the infinitesimal change
\[
\mu\to\mu+\delta\mu
\]
the renormalized field $[\varphi^2]_{(\mu)}$ transforms as
\[
[\varphi^2]_{(\mu+\delta\mu)}=\left(1-\frac{\delta\mu}\mu\,\gamma_{\varphi^2}(\lambda)\right)\,[\varphi^2]_{(\mu)}\,.
\]
This brings in certain change in the action, which however is compensated if we
simultaneously change the mass parameter,
\[
M^2\to\left(1+\frac{\delta\mu}\mu\,\gamma_{\varphi^2}(\lambda)\right)\,M^2\,.
\]
Invariance of the action w.r.t.\ the above simultaneous change in $\mu$ and $M^2$
leads to the equation
\[
\left(\mu\frac\partial{\partial\mu}+\beta(\lambda)\,\frac\partial{\partial\lambda}+n\,\gamma(\lambda)+\gamma_{\varphi^2}(\lambda)\,M^2\frac\partial{\partial M^2}\right)\bigl\langle\varphi(x_1)\cdots\varphi(x_n)\bigr\rangle=0\,,
\]
or, equivalently (using $\langle\varphi(x_1)\cdots\varphi(x_n)\rangle=\mu^n\,C^{(n)}(\mu x_i,M^2/\mu^2|\lambda)$),
\[
\left(\sum_ix_i\frac\partial{\partial x_i}+n\bigl(1+\gamma(\lambda)\bigr)+\beta(\lambda)\,\frac\partial{\partial\lambda}+\bigl(-2+\gamma_{\varphi^2}\bigr)\,M^2\frac\partial{\partial M^2}\right)\bigl\langle\varphi(x_1)\cdots\varphi(x_n)\bigr\rangle=0\,.
\]
This equation shows that in the massive case the scale transformations induce the
``flow'' in the two-dimensional space of $\varphi^4$ theories,
\begin{equation}\label{eq:massflow}
L\frac{\dd}{\dd L}\,\lambda=-\beta(\lambda)\,,\qquad L\frac{\dd}{\dd L}\,M^2=\bigl(2-\gamma_{\varphi^2}(\lambda)\bigr)\,M^2\,.
\end{equation}
Note that $\beta(\lambda)$ and $\gamma(\lambda)$ here are exactly the same as in the
massless theory, while $\gamma_{\varphi^2}(\lambda)$ is also determined from the
massless theory via Eq.~\eqref{eq:gammaphi2}.

\begin{hwproblem}{5}\label{hw:5}
In the massless $\varphi^4$ theory calculate $\Gamma^{(2)}$ and $\Gamma^{(4)}$
explicitly, to the order $\lambda_{(\mu)}^2$. Use dimensional regularization, and
the $\mu$-dependent scheme defined by Eqs.~\eqref{eq:mu1}, \eqref{eq:mu2}. In
particular, calculate the counterterm coefficients $z^{(2)}_{(\mu)}$ and
$a^{(2)}_{(\mu)}$, and confirm that $b^{(1)}_{(\mu)}=b^{(2)}_{(\mu)}=0$.

In solving this problem the following dimensionally continued integral is useful
\[
\int\frac{\dd^dk}{k^2}=0\,.
\]
\end{hwproblem}

\begin{hwproblem}{6}\label{hw:6}
Consider a scalar field theory with $N$-component field
$\bm{\varphi}=(\varphi_1,\dots,\varphi_N)$, and $O(N)$-symmetric action
\[
\mathcal A=\int \dd^dx\,\left[\frac12\,\bigl(\partial\bm{\varphi}\bigr)^2+\frac{M^2}2\,\bm{\varphi}^2+\frac\lambda8\,\bigl(\bm{\varphi}^2\bigr)^2\right]\,.
\]
\begin{hwparts}
\item Use the $O(N)$ symmetry to show that the two- and four-point correlation
functions $\tilde W^{(2)}_{ij}(p)$ and $\tilde W^{(4)}_{ijkl}(p_1,p_2,p_3,p_4)$
defined as
\begin{align*}
\tilde W^{(2)}_{ij}(p)&:=\text{f.t.}\,\bigl\langle\varphi_i(x)\varphi_j(x')\bigr\rangle\,,\\
\tilde W^{(4)}_{ijkl}(p_1,p_2,p_3,p_4)&:=\text{f.t.}\,\bigl\langle\varphi_i(x_1)\varphi_j(x_2)\varphi_k(x_3)\varphi_l(x_4)\bigr\rangle
\end{align*}
enjoy the following structures
\begin{align*}
\tilde W^{(2)}_{ij}(p)&=\delta_{ij}\,\tilde W^{(2)}(p)\,,\\
\tilde W^{(4)}_{ijkl}(p_1,p_2,p_3,p_4)&=\delta_{ij}\delta_{kl}\,\tilde W^{(4)}(p_1,p_2;p_3,p_4)+\delta_{ik}\delta_{jl}\,\tilde W^{(4)}(p_1,p_3;p_2,p_4)\\
&\qquad+\delta_{il}\delta_{jk}\,\tilde W^{(4)}(p_1,p_4;p_2,p_3)
\end{align*}
with scalar functions $\tilde W^{(2)}(p)$ and $\tilde W^{(4)}(p_1,p_2;p_3,p_4)$.
Define the scalar proper vertices $\Gamma^{(2)}(p)$ and
$\Gamma^{(4)}(p_1,p_2,p_3,p_4)$.
\item Introduce all elements of Feynman diagram technique (in the momentum space)
for this model, i.e.\ propagators and vertices.
\item For the massless theory write down the Callan-Symanzik equations for the
proper vertices $\Gamma^{(2)}(p)$ and $\Gamma^{(4)}(p_1,p_2,p_3,p_4)$.
\item At $d=4$ compute the functions $\beta(\lambda)$ and $\gamma(\lambda)$ to the
order $\lambda^2$.
\end{hwparts}
\end{hwproblem}
